\phantomsection
\addcontentsline{toc}{chapter}{ĐỀ CƯƠNG CHI TIẾT}

\begin{center}
{\large\bfseries{ĐỀ CƯƠNG CHI TIẾT}}
\end{center}

Đề cương chi tiết của khóa luận được tóm lược như sau:

\begin{enumerate}
    \item \textbf{Bối cảnh nghiên cứu:} sự bùng nổ của video ngắn làm gia tăng nhu cầu phân đoạn video theo thời gian một cách tự động và chính xác.
    \item \textbf{Bài toán đặt ra:} xác định ranh giới giữa các shot trong video, bao gồm cả chuyển cảnh đột ngột và chuyển cảnh dần dần.
    \item \textbf{Khó khăn chính:} nhịp cắt nhanh, nội dung biến đổi mạnh trong nội bộ shot, nhiều hiệu ứng biên tập và thiếu dữ liệu chuyên biệt cho video ngắn.
    \item \textbf{Dữ liệu nghiên cứu:} bộ dữ liệu SHOT được dùng làm tập trung tâm, đồng thời tham chiếu thêm các bộ dữ liệu BBC Planet Earth và ClipShots để đối chiếu.
    \item \textbf{Phương pháp:} khảo sát TransNetV2 và AutoShot; phân tích cơ chế NAS; triển khai hướng cải tiến AutoShotV2 dựa trên tầng phân loại mới, hàm mất mát Focal, nhãn nhiều điểm dương, hiệu chỉnh nhiệt độ và làm mượt Gaussian.
    \item \textbf{Thực nghiệm:} đánh giá bằng các chỉ số độ chính xác dương, độ bao phủ và chỉ số F1; so sánh mô hình cơ sở TransNetV2, AutoShot gốc và AutoShotV2 trên tập SHOT.
    \item \textbf{Kết quả chính:} Ở ngưỡng cho F1 cao nhất trên từng bộ dữ liệu, AutoShotV2 đạt F1 = \ASVTwoBestShotFOne{} trên tập kiểm tra SHOT, cao hơn AutoShot tự chạy lại ($0.8405$) và vượt đáng kể mô hình cơ sở TransNetV2 ($0.7993$); đồng thời vẫn tổng quát hóa tốt trên BBC và ClipShots.
    \item \textbf{Kết cấu báo cáo:} gồm năm chương, lần lượt trình bày tổng quan, công trình liên quan, phương pháp, dữ liệu và thực nghiệm, kết luận và hướng phát triển.
\end{enumerate}
